\documentclass[12pt,a4paper]{article}
\usepackage[utf8]{inputenc}
\usepackage[spanish]{babel}
\usepackage{geometry}
\usepackage{parskip}
\geometry{a4paper, margin=2.5cm}

\title{\textbf{Tarea 1}\\ \large Bases de datos grupo 1}
\author{Angel Gabriel Vazquez Toledo}
\date{} % Omitimos la fecha o la puedes agregar con \today

\begin{document}

\maketitle

\section*{Características de la información}
\begin{itemize}
    \item \textbf{Veraz:} Debe de provenir de fuentes confiables. No tener errores y manipulaciones, debe de estar verificada.
    \item \textbf{Relevante:} Debe de ser útil para el usuario y estar alineada con el propósito que se quiere cumplir, cumpliendo objetivos e intereses.
    \item \textbf{Actual:} Debe de ser reciente y vigente, no debe de ser obsoleta.
    \item \textbf{Accesible y disponible:} La información debe de ser fácil de localizar, obtener y ser localizada por el usuario.
    \item \textbf{Confidencialidad:} Los datos deben de estar protegidos ante usos no autorizados o indebidos, como por ejemplo la información sensible de un usuario como lo son sus datos personales.
    \item \textbf{Integridad:} La información debe ser completa, consistente y coherente. No debe de presentar pérdida de información.
\end{itemize}

\section*{Modelos de datos}

\textbf{Modelo orientado a objetos:} \\
El modelo orientado a objetos guarda la información directamente como objetos, incluyendo sus atributos y métodos. Su gran ventaja es que evita tener que traducir el código a tablas relacionales, agilizando el desarrollo. Sin embargo, su mantenimiento es complejo, por lo que su uso actual se limita como simulaciones científicas o software de diseño.

\textbf{Modelo columnar:} \\
El modelo columnar organiza los datos físicamente por columnas. Esto lo hace rapidísimo para consultas analíticas masivas, ya que el motor solo lee la columna que necesitas evaluar. Su desventaja es que resulta ineficiente para insertar o actualizar registros uno por uno. Es el estándar para Big Data y Data Warehousing.

\textbf{Modelo documental:} \\
El modelo documental almacena la información en formatos similares a JSON sin requerir un esquema fijo. Esta flexibilidad es su mayor ventaja, permitiendo tener estructuras diferentes en una misma colección. Su desventaja es que hace cruces de información (joins) complejos y consume muchos recursos. Se utiliza mucho en aplicaciones web modernas y catálogos.

\textbf{Modelo clave-valor:} \\
El modelo clave-valor funciona como un diccionario: un identificador único apunta directamente a su contenido. Es increíblemente rápido para leer y escribir, pero tiene la desventaja de que no permite hacer consultas sobre el contenido, solo puedes buscar con la clave exacta. Es ideal para sistemas de caché o manejo de sesiones temporales.

\textbf{Modelo orientado a grafos:} \\
El modelo orientado a grafos se basa en nodos (entidades) y aristas (relaciones). Su principal ventaja es que puede analizar y recorrer millones de conexiones en milisegundos, algo que saturaría a una base relacional tradicional. Como desventaja, es ineficiente si los datos son planos y poco conectados. Su fuerte son las redes sociales, sistemas de recomendación y detección de fraudes.

\end{document}
